% STATUS: written 2026-11-05 (Singer-Wermer 1955 read in full; excerpts in
% excerpts/SingerWermer1955.txt; Thomas 1988 read from the library copy
% 2026-11-09 and cited first-hand).
% Additions of 2026-11-07 from Mathieu 1994 (pp. 238, 242--245) and
% Eisaei--Moghimi 2024 (pp. 1--5).
% Sources: SingerWermer1955, BresarVillena2002, Mathieu1994Image,
% EisaeiMoghimi2024Theta.
\section{Singer--Wermer and commutative Banach algebras}\label{sec:banach-operator:singer-wermer}

The chapter opened with the observation of \v{S}ilov that the algebra of
infinitely differentiable functions on an interval carries no Banach
algebra norm \cite[p.~260]{SingerWermer1955}. Singer and Wermer took up
the explanation of Kaplansky of this phenomenon and worked with the following
objects \cite[p.~260]{SingerWermer1955}.

\begin{definition}[{\cite[p.~260]{SingerWermer1955}}]\label{def:banach:derivation}
A \emph{derivation} $D$ on an algebra is a transformation on the algebra
such that $D(a+b)=D(a)+D(b)$, $D(\lambda a)=\lambda D(a)$ for any scalar
$\lambda$, and $D(ab)=D(a)b+aD(b)$. A derivation on a Banach algebra is
\emph{bounded} if $\sup_{\|a\|=1}\|D(a)\|<\infty$. The \emph{radical} of a
commutative Banach algebra $A$ is the intersection of all maximal ideals
$M$ in $A$ such that $A/M$ has a unit; if the radical reduces to the zero
element, $A$ is called \emph{semisimple}.
\end{definition}

Their main theorem answers the conjecture of Kaplansky in the bounded case
\cite[p.~260]{SingerWermer1955}.

\begin{theorem}[{\cite[Theorem~1, p.~260]{SingerWermer1955}}]\label{thm:banach:singerwermer}
Let $A$ be a commutative Banach algebra and $D$ a bounded derivation on
$A$. Then $D$ maps $A$ into its radical. In particular, if $A$ is
semisimple, $D=0$.
\end{theorem}

The proof of Theorem~\ref{thm:banach:singerwermer} uses the exponential.
For a multiplicative linear functional $f$ and a complex number $\lambda$,
the map $a\mapsto f(e^{\lambda D}a)$ is again multiplicative. Gelfand
bound on multiplicative functionals then forces $f(D(a))$ to vanish
\cite[p.~260]{SingerWermer1955}. The result of \v{S}ilov is its first
consequence \cite[p.~260]{SingerWermer1955}.

\begin{corollary}[{\cite[Corollary~1.2, p.~262]{SingerWermer1955}}]\label{cor:banach:cinfty}
Let $C^\infty$ denote the algebra of all infinitely differentiable
complex-valued functions on an interval. Then there exists no norm under
which $C^\infty$ is a Banach algebra.
\end{corollary}

For the proof, differentiation $f\mapsto f'$ would be a bounded derivation
by the closed graph theorem, and $C^\infty$ is semisimple, so
Theorem~\ref{thm:banach:singerwermer} would give $D=0$
\cite[p.~262]{SingerWermer1955}. Singer and Wermer added two observations
that fed the later commutator theory of
Section~\ref{sec:banach-operator:commutator-ranges}. A bounded commutator
$ab-ba$ inside a commutative Banach algebra is a generalized nilpotent
\cite[p.~261]{SingerWermer1955}. The theorem of Putnam, that a commutator
$C=AB-BA$ of bounded Hilbert space operators commuting with both $A$ and
$B$ has spectrum $\{0\}$, follows from Theorem~\ref{thm:banach:singerwermer}
\cite[p.~264]{SingerWermer1955}.

Derivations into larger algebras behave differently, and the obstruction
lives at one point of the maximal ideal space
\cite[p.~262]{SingerWermer1955}.

\begin{definition}[{\cite[p.~262]{SingerWermer1955}}]\label{def:banach:pointderivation}
A \emph{bounded point derivation} of a commutative Banach algebra $A$ is a
bounded linear functional $d_\varphi$ associated with a multiplicative
linear functional $\varphi$ such that
$d_\varphi(a_1a_2)=d_\varphi(a_1)\varphi(a_2)+\varphi(a_1)d_\varphi(a_2)$.
\end{definition}

The next theorem relates point derivations to derivations into a larger algebra.

\begin{theorem}[{\cite[Theorem~2, p.~262]{SingerWermer1955}}]\label{thm:banach:pointder}
If there exists a nonzero bounded point derivation $d_\varphi$ of $A$, then
there exists a commutative extension $B$ of $A$ and a nonzero bounded
derivation $D$ of $A$ into $B$. If $A$ is semisimple, $B$ can be taken to
be semisimple. If $A\subseteq B$ and if $D$ is a nonzero bounded derivation
of $A$ into $B$ but not into the radical of $B$, then there exists a
nonzero bounded point derivation of $A$.
\end{theorem}

On a compact Hausdorff space the point derivation option disappears
\cite[p.~264]{SingerWermer1955}.

\begin{corollary}[{\cite[Corollary~2.2, p.~264]{SingerWermer1955}}]\label{cor:banach:CX}
The algebra $C(X)$ of all continuous functions on a compact Hausdorff space
$X$ has no nonzero derivations into any semisimple commutative extension
$B$.
\end{corollary}

Singer and Wermer wrote that the boundedness hypothesis of
Theorem~\ref{thm:banach:singerwermer} seemed superfluous
\cite[p.~260]{SingerWermer1955}. The statement became the
Singer--Wermer conjecture \cite[p.~523]{MiuraHirasawaTakahasi2006}.
Thomas proved it in 1988: a
linear derivation on a commutative Banach algebra maps into the radical
even if it is discontinuous \cite[pp.~435--460]{Thomas1988Image}. As a
direct consequence, there are no nonzero linear derivations on a semisimple
commutative Banach algebra, a fact that Johnson had proved earlier
\cite[p.~523]{MiuraHirasawaTakahasi2006}. That consequence is the
commutative end of the continuity story told in
Section~\ref{sec:banach-operator:automatic-continuity}. In the
noncommutative algebras of
Section~\ref{sec:banach-operator:kadison-sakai}, Singer showed that a commutative
$C^*$-algebra has only the zero derivation
\cite[pp.~260--261]{SingerWermer1955}. Kaplansky proved every
derivation of a type I von Neumann algebra inner
\cite[p.~280]{Kadison1966}.

The noncommutative form of the conjecture is open, and this section states it
as a problem. % claim-ok: own signpost introducing the moved problem, which carries its own citation

\begin{problem}[{\cite[p.~710]{BresarVillena2002}}]\label{prob:open:sw}
Let $A$ be a Banach algebra and let $D$ be a derivation of $A$, with no
continuity assumed.
Does $D$ leave every primitive ideal of $A$ invariant?
\end{problem}

Bre\v{s}ar and Villena call this the noncommutative Singer--Wermer conjecture
and describe it as one of the hardest open problems of Banach algebra theory
\cite[p.~710]{BresarVillena2002}. They prove that it is equivalent to a
statement about invertible values of a derivation
\cite[Theorem~5.2, p.~719]{BresarVillena2002}. They also report that the second
assertion of their Theorem~5.2 is known for derivations but not for
$\varphi$-derivations \cite[Remark~5.3, p.~720]{BresarVillena2002}.

Mathieu surveyed the state of the conjecture in 1994 and put it in two forms
\cite[p.~238]{Mathieu1994Image}. A derivation is \emph{commuting} when
$[a,\delta(a)]=0$ for every element $a$ \cite[p.~238]{Mathieu1994Image}. The
conjecture then reads as follows: every derivation that is commuting modulo the
radical is inner modulo the radical \cite[Conjecture~1, p.~238]{Mathieu1994Image}.
The second form is a statement about the image of one element
\cite[p.~242]{Mathieu1994Image}.

\begin{conjecture}[{\cite[Conjecture~2, p.~242]{Mathieu1994Image}}]\label{conj:banach:unbounded-kleinecke}
Let $\delta$ be a derivation on a Banach algebra $A$. For each $a$ in $A$,
$[a,\delta(a)]=0$ implies that $\delta(a)$ is quasinilpotent.
\end{conjecture}

Mathieu calls Conjecture~\ref{conj:banach:unbounded-kleinecke} the unbounded
Kleinecke--Shirokov theorem and proves that it follows from the
conjecture on primitive ideals \cite[p.~242]{Mathieu1994Image}. Three results of that survey bound the image from the other side
\cite[pp.~243--245]{Mathieu1994Image}. A derivation
that maps into the centre maps into the radical
\cite[Theorem~7, p.~243]{Mathieu1994Image}. No derivation of a unital Banach
algebra reaches an invertible central value
\cite[Theorem~8, p.~245]{Mathieu1994Image}. In particular, no derivation of a
unital Banach algebra takes the value $1$
\cite[Corollary~9, p.~245]{Mathieu1994Image}. Mathieu notes that the last
statement was the key step in the proof of Thomas of the commutative conjecture
\cite[p.~245]{Mathieu1994Image}.

The conjecture also has a form for twisted product rules
\cite[p.~1]{EisaeiMoghimi2024Theta}. Eisaei and Moghimi work with a
weighted space of functions on the half-line, and with its dual
\cite[p.~2]{EisaeiMoghimi2024Theta}. That dual is a Banach algebra under the
first Arens product, and the paper writes it $L^\infty_0(w)^*$
\cite[p.~2]{EisaeiMoghimi2024Theta}. They twist the rule by a homomorphism
$\theta$ of that algebra. A linear map $D$ is a
\emph{$\theta$-derivation} when $D(mn)=D(m)\theta(n)+\theta(m)D(n)$ for all
$m,n$ \cite[p.~1]{EisaeiMoghimi2024Theta}. The radical of that algebra equals
its right annihilator \cite[p.~3]{EisaeiMoghimi2024Theta}. For an isomorphism
$\theta$ the range of $D$ lies in the radical
\cite[Theorem~2.1, p.~4]{EisaeiMoghimi2024Theta}. If $\theta$ is also
continuous, then $D$ is continuous exactly when its restriction to the right
annihilator is continuous \cite[Proposition~2.3, p.~5]{EisaeiMoghimi2024Theta}.

Eisaei, Mehdipour and Moghimi treat the one-sided rule in the same algebra
\cite[p.~2]{EisaeiMehdipourMoghimi2024LeftTheta}. A linear map $\delta$ is a
\emph{left $\theta$-derivation} when
$\delta(st)=\theta(s)\delta(t)+\theta(t)\delta(s)$ for all $s,t$
\cite[p.~2]{EisaeiMehdipourMoghimi2024LeftTheta}. Such a map is already a
$\theta$-derivation, so the one-sided rule adds nothing in this algebra
\cite[Theorem~2.2, p.~4]{EisaeiMehdipourMoghimi2024LeftTheta}. Continuity
behaves as before: if a right identity $\nu$ satisfies $\|\theta(\nu)\|<1$,
then the zero map is the only continuous left $\theta$-derivation
\cite[Corollary~2.5, p.~6]{EisaeiMehdipourMoghimi2024LeftTheta}. When $\theta$ is injective, the range of
$\delta$ lies in the radical exactly when $\delta=0$
\cite[Theorem~2.7, p.~7]{EisaeiMehdipourMoghimi2024LeftTheta}.

Russell treats the automorphisms and derivations of continuous trace
$C^{*}$-algebras together, which is where the inner derivations of such an
algebra are classified \cite[Abstract, p.~1]{Russell1980ContinuousTrace}.

Standard operator Jordan algebras on a Banach space are introduced and their derivations determined, extending the standard operator algebra case of the same chapter~\cite{WlStandardOperatorJordan2022}.

The noncommutative Singer-Wermer conjecture, that every linear derivation of a Banach algebra maps into its Jacobson radical, is studied through linear left derivations~\cite{WlNoncommutativeVersionsOf2008}.

The noncommutative Singer-Wermer conjecture, that every linear derivation of a Banach algebra maps into its Jacobson radical, is studied through linear left derivations~\cite{WlNoncommutativeVersionsOf2008}.

Unbounded reflexive $*$-derivations of $C^*$-algebras of bounded operators are studied with weakly dense domains containing the compact operators, and their lattices of invariant subspaces are described~\cite{WlReflexiveStarDerivationsAndLat2017}.

For certain classes of Banach algebras the only derivation is zero, which is the vanishing form of the Singer-Wermer theorem the chapter states~\cite{WlNonexistenceOfNonzeroDerivatio2014}.

Moumen, Taoufiq and Boua prove that if a non-injective continuous derivation of a Banach algebra over $\mathbb{R}$ or $\mathbb{C}$ satisfies certain local differential identities, then the algebra is commutative~\cite{WlOnPrimeBanachAlgebrasWithConti2023}.

Hosseini shows that if the self-adjoint elements of a Banach $*$-algebra $A$ are real-linear combinations of mutually orthogonal projections and a uniformly bounded sequence of linear maps satisfies the stated Leibniz-type relation at every projection, then the maps take their values in the intersection of the kernels of the characters of $A$, and vanish when $A$ is semiprime and that intersection has dimension at most one~\cite{WlSomeConditionsUnderWhichDeriva2017}.

Ouzomgi studies point derivations of finite and infinite order at a character of a commutative Banach algebra, the functionals $d_k$ with $d_k(fg)=\sum_{j=0}^{k}d_j(f)d_{k-j}(g)$, and relates the ideals they define to the structure of Banach algebras of differentiable functions~\cite{WlPointDerivationsAndIdealsInBan1989}.

Lyubich and Phong prove that if $A$ generates a strongly continuous bounded semigroup, the intersection of the spectrum of $A$ with the imaginary axis is at most countable, and the adjoint $A^*$ has no purely imaginary eigenvalues, then the Cauchy problem for $\dot{x}(t)=Ax(t)$ is asymptotically stable~\cite{WlAsymptoticStabilityOfLinearDif1988}.

Read exhibits a class of Banach spaces $E$ for which discontinuous derivations of the algebra $B(E)$ of bounded operators exist, the first examples of their kind, and finds new Banach spaces in the process~\cite{WlDiscontinuousDerivationsOnTheA1989}.
